\documentclass[literature.tex]{subfiles}

\begin{document}
\prose{Obraz Doriana Graye}{Oscar Wilde}
\teorieA{
Román, obsahuje hororové prvky
}{}{}{}{}{
Text je rozdělen do dvaceti kapitol.
Paralérní kompozice děje, kdy se takto seznamujeme se Sibyl Vaneovou a její rodinou.
Autor na začátek ještě přidal \textbl{předmluvu}, která slouží jako jakýsi \texthl{návod na pochopení jeho díla}.
}{
\enquote{Ty nechceš, abych se s ním setkal?} \\
\enquote{Ne.} \\
\enquote{Pan Dorian Gray je v ateliéru, pane,} pravil komorník, vstupující do zahrady. \\
\enquote{Teď mě představit musíš,} řekl lord Henry se smíchem. \\
Malíř se obrátil k sluhovi, který tu stál a pomrkával ve slunečním světle.
\enquote{Požádejte pana Graye, aby počkal, Parkere. Že přijdu za~okamžik.}
Sluha se uklonil a vracel se po pěšině.
Basil se zahleděl na lorda Henryho.
\enquote{Dorian Gray je můj nejdražší přítel,} řekl.
\enquote{Má prostou a krásnou povahu.
Tvoje teta ti o něm řekla pravdu.
\texthl{Nezkaz ho.
Nesnaž se ho ovlivnit.
Tvůj vliv by byl špatný.}
Svět~je~veliký a je v něm plno skvělých lidí.
Neodveď mi tedy jedinou bytost, které moje umění vděčí za veškeré kouzlo, co v něm je.
Závisí na něm můj život umělce.
Na to nezapomeň, Harry, důvěřuji ti.}
Mluvil velmi zvolna a zdálo se, že se ta slova z něho derou téměř proti jeho vůli. \\
\enquote{Co to říkáš za nesmysl!} řekl lord Henry s úsměvem, a uchopiv Hallwarda pod paží,
téměř ho vlekl do domu. 
}
\teorieB{
\wtem{Dorian Gray}{hlavní postava, jehož příběh sledujeme. Začíná jako nádherný a nezkažený mladý muž, ještě adolescent. Díky vlivům \textbl{svých, Basilových a Lorda Henryho} se ale postupně kazí,
až spáchá dokonce i \texthl{vraždu}. Už od útlého věku se u něj projevuje narcisimus a během svého života se propadá stále hloubš \wemph{za honbou po potešení}. Stává se z něj proutník, narkoman a jeho osobnost upadá.}
\wtem{Basil Harvard}{idealistický malíř, Graye si všimne jako jeden z prvních a \texthl{zamiluje se do něj}. I jeho nekonečné velebení Dorianovi osobnosti mohlo mít vliv na jeho budoucnost.}
\wtem{Lord Henry}{manipulátor, miluje ovlivňovat okolí a sledovat jak se vyvine. To Lord Henry byl první, který Dorianovi nabídl svět potěšení a i když by si ho narcistický Dorian dost možná našel i sám, tak mu to nikdy nepřestal vyčítat. Během díla vystupuje jako filozof a inteligentní muž.}
\wtem{Sibyl Vaneová}{první dívka, jež měla tu smůlu a zamilovala se do Doriana Graye. \wemph{Také první život, který kdy Dorian vzal} a který byl předzvěstí jeho dalšího počínání.}
}{
Příběh začíná pohledem na místo Basila Harvarda, ke kterému na návštěvu přišel jeho přítel Lord Henry.
Baví se o mladíkovi jménem Dorian Gray
a~tím uvádějí jeho charakter ještě před uvedením jeho postavy do děje.
Po této první kapitole už se děj odehrává výhradně kolem Doriana Graye.
Zamiluje se do Sibyl Vaneové, jen aby podlehl potěšení a nepřímo zavinil její smrt.
Po této zkušenosti propadne požitkářství bohémského života a~postupně se propadá.
Tráví noci ve vykřičených domech a opiových doupatech, okolí se od něj distancuje
a to celé zatímco na něm se \texthl{nezmění ani vráska}.
Jeho~modlitba vyřčena toho prvního dne u Basila Harvarda totiž zavinila,
že všechny projevy stáří i morálního úpadku \texthl{mění podobu jeho obrazu}
a~nikoliv podobu Doriana samotného.
Ten, co Doriana neviděl něco z toho dělat ani neuvěří,
že ten mile vypadající adolescent je třicetiletý proutník závislý na opiu.
Během tohoto morálního úpadku \texthl{zabije Basila Harvarda},
který namaloval obraz a v jeho očích to celé zavinil.
I Lordu Henrymu vyčítá svojí situaci.
Ale děj nikdy neurčil, \wemph{kdo ze tří mužů mohl za situaci Doriana Graye},
zda on sám, jeho přátelé, či všichni dohromady.
A příběh končí, když se Dorian rozhodne \texthl{zničit obraz},
důkaz o jeho morálním úpadku a tím \textbl{zabije sám sebe}.
}{
Většina děje, až na pár výjimek, se odehrává v Londýně. Doba je circa konec 19. století.
}{
Autor zde pojednává o muži, jenž se behlavě žene za potěšením a tím nezničí jen sám sebe, ale i lidi kolem sebe.
I když hovoří o umění a kráse potěšení, tak mluví i o trestu, který má být očistou a kterému se Dorian ve své nepřirozenosti vyhýbá.
}
\historie{}{}{}{}{}{
Oscar Wilde byl jistě geniální muž, který rád dráždil své okolí nemístníma narážkama.
Rád nechával své okolí v nejistotě, zda není homosexuál a to ho přišlo draho, když byl odsouzen za homosexualitu ke dvěma rokům těžkých prací.
Po odsloužení trestu se už nikdy nedal zcela dohromady a umřel v~bídě jen pár let poté.
}{
Mezi nejznámnější díla patří \textsc{Jak je důležité míti Filipa}, právě \textsc{Obraz Doriana Graye}, \textsc{Strašidlo cantervillské}, a \textsc{Lidská duše za socialismu}.
}\newpage
\kritika{
Když Oscar Wilde původně vydal Obraz Doriana Graye do novin, tak to jeho nakladateli přišlo až moc pohoršující a pozměnil či smazal asi 500 slov.
Tento muž, Stoddart, mazal jak homosexuální, tak i heterosexuální narážky.
Gray například o svých ženách dřív mluvil jako o milenkách, ale i jeho vztah s Basilem Harvardem byl hlubší.\par
Tato cenzura však nebyla konečná.
Dílo se stejně setkalo s kritikou a sám Oscar Wilde, když se dílo rozhodl vydat jako knihu, tak provedl autocenzuru.
Pozměnil některé texty, přidal další tři kapitoly. Poslední kapitolu rozdělil na dvě. Přidal i předmluvu, která je jeho odpovědí všem kritikům původní, časopisecké, verze.\par
Stejně byl odsouzen ke dvěma letům nucených prací a právě toto dílo bylo použito jako jeden z důkazů, že má homosexuální sklony.
}{}{
Bohémský život se vyskytuje stále a i když už se homosexualita nevnímá jako problém, tak přehnané požitkářství ano.
Oscar Wilde představuje osobu Doriana Graye bravurně a je až hrůzostrašné sledovat, jak jednoduše se dají zničit životy lidí kolem sebe, když je člověk jen trochu neohleduplný.
Vlastně se člověk může do Doriana Graye poměrně snadno vcítit a to to horší dopad na čtenáře má, když najednou vidí dopady jeho chování.
}
\nazor{
Wildeovo dílo mě chytilo a až do konce mě nepustilo.
Nabídlo mi nový pohled na svět kolem sebe, i lekci, kterou by si měli poslechnout všichni.
Naprosto rozumím, proč je Obraz Doriana Graye jednou z nejdůležitějších literatur své doby.
}
%\explain{}
\wpage
\end{document}